%!TEX root = ../AdvMathLecture.tex
%
% Section: Introduction
%

%~~~~~~~~~~~~~~~~~~~~~~   []   ~~~~~~~~~~~~~~~~~~~~~~~~~~~
\begin{frame}
\frametitle{\LARGE {\color[rgb]{0.02,0.20,1.0} \UL{Introduction}}}

\vspace{0.2cm}

\begin{columns}
\begin{column}{5.0cm}
	I. Functions
\end{column}
\begin{column}{6.8cm}	\end{column}
\end{columns}

\scriptsize
\begin{columns}
\begin{column}{12.0cm}


\begin{columns}
\begin{column}{7.5cm}
	%-------LIST---------
\begin{itemize}
	\item Functions have \ON<2->{\B{inputs and outputs.}}
	\item Inputs live in the \ON<2->{\B{Domain.}}
	\item Outputs live in the \ON<2->{\B{Range.}}
	\item For each Input there can be \ON<2->{\B{ONLY ONE Output.}}
	\item 2 different inputs can correspond to \ON<2->{\B{1 Output.}}
	\item 1 input can \ON<2->{\B{NEVER correspond to 2 different Outputs.}}
\end{itemize}
%-------------------
\end{column}
\begin{column}{4.0cm}
%	The point of a function is to take ``objects" in the Domain and then do ``something" to them. The new objects now live in the Range.   
	
\centering

$\D{F(Domain) \longrightarrow Range}$

\vspace{0.2cm}
	
\hspace{0.2cm} $\D{F(Inputs) \hspace{0.1cm} \longrightarrow Outputs}$

\end{column}
\end{columns}

\vspace{0.3cm}

In Calculus 1 and 2, we had $f(x)$ ...

%-------LIST---------
\begin{itemize}
	\item Domain/Input: \quad $x \in \RE$ \, (1D line, real numbers)
	\item Range/Output: \quad $f(x) \in \RE$
\end{itemize}
%-------------------

\vspace{0.1cm}

In Calculus 3, we had $f(x,y)$ and $f(x,y,z)$ ... 

%-------LIST---------
\begin{itemize}
	\item Domain/Input: \quad $(x,y) \in \RRE$ \, (2D plane) \, or \, $(x,y,z) \in \RRRE$ \, (3D space)
	\item Range/Output: \quad $f(x,y) \in \RE$ \, or \, $f(x,y,z) \in \RE$
\end{itemize}
%-------------------

\vspace{0.1cm}

In This Class ... we consider $\VF(x)$, $\VF(x,y)$, and $\VF(x,y,z)$ ... 

%-------LIST---------
\begin{itemize}
	\item Domain/Input: \quad $x \in \RRE$ (2D plane) \, or \, $\RRRE$ (3D space)
	\item Range/Output: \quad $\VF(x) \in \RE$
\end{itemize}
%-------------------

\end{column}
\begin{column}{0.0cm}\end{column}
\end{columns}
\end{frame}

%~~~~~~~~~~~~~~~~~~~~~~~~~~~~~~~~~~~~~~~~~~~~~~~~~~~~~~~~~
\begin{frame}
\frametitle{}
\scriptsize
\begin{columns}
\begin{column}{12.0cm}
	
	\DEF A {\bf vector}, ${\bf v}$, is a mathematical object that has both \ON<2->{\B{\quad magnitude/length \quad }} and \ON<3->{\B{\quad direction.}}

\vspace{0.3cm}

\UL{Notation}: \\

\vspace{0.2cm}

\OBUL  Vectors will written either in bold or with an arrow: \quad {\bf v}, $\vec{v}$.\\

\vspace{0.2cm}

\OBUL  A vector, $\vec{v}$, in $\mathbf{R}^2$ (\,$\vec{v} \in \mathbf{R}^2$\,) can be represented by \ON<6->{\B{\,2\,}} coordinates: \quad $\vec{v} = \ON<6->{\B{\VV \in \mathbf{R}^2}}$\\

\vspace{0.2cm}

\OBUL  A vector, $\vec{v}$, in $\mathbf{R}^3$ (\,$\vec{v} \in \mathbf{R}^3$\,) can be represented by \ON<7->{\B{\,3\,}} coordinates: \quad $\vec{v} = \ON<7->{\B{\VVV \in \mathbf{R}^3}}$\\

\vspace{0.2cm}

\OBUL  The special unit vectors {\bf i} = [1,0,0], {\bf j} = [0,1,0], and {\bf k} = [0,0,1] form a basis for $\RE^3$, \NEWLINE so we can write \,\, $\vec{v} = \VVV$ \,\, as \,\, \ON<8->{\B{$\vec{v} = v_1{\bf i} + v_2{\bf j} + v_3{\bf k}$ \quad (linear combination of \VI, \VJ, \VK)}}

\vspace{0.3cm}

\EX \\
\vspace{0.1cm}
\OBUL $[10,14], [0,1], [-1,2] \ON<9->{\B{\,\, \in \mathbf{R}^2}}$ \quad and \quad $[10,4] = \ON<10->{\B{10\VI + 4 \VJ}}$\\
\vspace{0.2cm}
\OBUL $[1,1,-1], [0,0,0], [-10,12,-8] \ON<11->{\B{\,\, \in \mathbf{R}^3}}$ \quad and \quad $[-10,12,-8] = \ON<12->{\B{-10\VI + 12 \VJ - 8\VK}}$\\

\vspace{0.5cm}

\THM Let ${\bf a} = [a_1,a_2,a_3]$ and ${\bf b} = [b_1,b_2,b_3]$ be vectors. Then
\begin{center}
	{\bf a} = {\bf b} \,\, $iff$ \quad \ON<13->{\B{$a_1 = b_1$, $a_2 = b_2$, $a_3 = b_3$}}
\end{center}
\EX \quad $[2,0,2d+c] = [c,0,1] \quad \implies$ \quad \ON<14->{\B{$c = 2$ \quad and \quad $2d+c=1$,}} \quad \ON<15->{\B{so \quad $d=-\frac{1}{2}$}}

%-------------------

\end{column}
\begin{column}{0.0cm}\end{column}
\end{columns}
\end{frame}

%~~~~~~~~~~~~~~~~~~~~~~~~~~~~~~~~~~~~~~~~~~~~~~~~~~~~~~~~~
\begin{frame}

\begin{columns}
\begin{column}{5.0cm}
	II. Vectors \& Vector Fields
\end{column}
\begin{column}{6.8cm}	\end{column}
\end{columns}

\scriptsize
\begin{columns}
\begin{column}{12.0cm}
\UL{Vector Operation 1}: Addition

\vspace{0.2cm}

Let \,\, $\VA = \AAA$ \,\, and \,\, $\VB = \BBB$

\vspace{0.2cm}

Then 
%--------EQN----------
\begin{align*}
	\VA + \VB = \AAA + \BBB = \ON<2->{\B{[a_1+b_1, a_2+b_2, a_3+b_3]}}
\end{align*}
%----------------------

\begin{center}
	\only<1-2>{\includegraphics[scale = 0.19]{Figures/9_1_Fig1.png}}
	\only<3->{\includegraphics[scale = 0.19]{Figures/9_1_Fig2.png}}
	\ON<4->{\B{\includegraphics[scale = 0.19]{Figures/9_1_Fig3.png}}}
\end{center}

\UL{Note}:

\vspace{0.1cm}

\OBUL $\VA,\VB \in \RE^3 \quad \implies \quad \ON<5->{\B{\VA + \VB \in \RE^3}}$ \\

\vspace{0.1cm}

\OBUL $\VA + \VB$ is known as the \ON<6->{\B{\,\,resultant\,\,}} of $\VA$ and $\VB$. (e.g. If $\VA$ and $\VB$ were forces then $\VA + \VB$ would be \NEWLINE the \ON<7->{\B{resultant force.}}
\end{column}
\begin{column}{0.0cm}\end{column}
\end{columns}

\vspace{10.2cm}

\end{frame}

%~~~~~~~~~~~~~~~~~~~~~~~~~~~~~~~~~~~~~~~~~~~~~~~~~~~~~~~~~
\begin{frame}

\begin{columns}
\begin{column}{5.0cm}
	III. Integrals
\end{column}
\begin{column}{6.8cm}	\end{column}
\end{columns}


\scriptsize
\begin{columns}
\begin{column}{12.0cm}
	
\end{column}
\begin{column}{0.0cm}\end{column}
\end{columns}

\vspace{10.2cm}

\end{frame}